\subsection{멀티 에이전트 구성 및 역할 정의}
본 연구에서는 복잡한 문서 생성 작업을 체계적으로 수행하기 위해, 전문화된 역할을 가진 네 가지 에이전트로 구성된 멀티 에이전트 시스템(Multi-Agent System, MAS)을 설계하였다. 각 에이전트는 특정 단계의 과업에 최적화된 프롬프트 엔지니어링을 통해 제어되며, 그 역할은 다음과 같다.

첫째, 설계 에이전트(Planner)는 입력된 요구사항과 기초 데이터를 분석하여 문서의 전체 구조와 목차를 설계한다. 이는 최종 결과물의 논리적 일관성을 결정하는 단계로, 각 섹션에 포함되어야 할 핵심 내용과 기술적 요구사항을 정의한 설계도를 생성한다. 둘째, 집필 에이전트(Writer)는 설계 에이전트가 작성한 설계도를 바탕으로 실제 본문을 작성한다. 이때 제공된 정형 데이터나 참조 자료를 활용하여 상세 내용을 구체화하며, 지정된 톤 앤 매너를 유지하는 데 집중한다. 셋째, 검수 에이전트(Reviewer)는 생성된 문서의 정확성, 완결성 및 설계도 준수 여부를 평가한다. 검수 에이전트는 단순한 오타 수정을 넘어, 논리적 비약이나 데이터 오류를 식별하여 구체적인 수정 피드백을 생성한다. 마지막으로, 수정 에이전트(Editor)는 검수 에이전트의 피드백을 반영하여 문서를 보완한다. 수정 에이전트는 원문의 의도를 유지하면서 검수 단계에서 지적된 결함을 제거하여 문서의 품질을 최종적으로 향상시키는 역할을 수행한다.

\subsection{오케스트레이션 구조별 제어 로직}
에이전트 간의 협업 흐름을 제어하는 오케스트레이션 방식에 따라 시스템의 동작 메커니즘은 결정론적 방식과 동적 방식으로 구분된다. 세 가지 구조의 상세 설계는 그림~\ref{fig:arch_compare}에 제시되어 있으며, 각 로직의 상세 내용은 다음과 같다.

\subsubsection{Code 기반 오케스트레이션}
Code 구조는 개발자가 사전에 정의한 상태 전이 그래프(State Transition Graph)에 따라 흐름이 제어되는 결정론적(Deterministic) 방식이다. 제어 로직은 하드코딩된 워크플로우 엔진에 의해 관리되며, 상태 전이 함수 $f: S_t \rightarrow S_{t+1}$는 고정된 규칙에 따라 작동한다. 예를 들어, $\text{Planner} \rightarrow \text{Writer} \rightarrow \text{Reviewer} \rightarrow \text{Editor} \rightarrow \text{End}$ 순으로 순차적으로 진행된다. 만약 검수 단계에서 반려(Reject) 판정이 날 경우, 미리 정의된 횟수($N$)만큼만 수정 단계로 되돌아가는 루프를 수행한다. 이 구조는 실행 경로가 예측 가능하여 자원 소비가 일정하며, 오버헤드가 최소화된다는 특징이 있다.

\subsubsection{LLM 기반 오케스트레이션}
LLM 구조는 별도의 라우터 LLM(Router LLM)이 현재의 작업 상태와 에이전트의 출력값을 실시간으로 분석하여 다음 경로를 결정하는 동적 라우팅(Dynamic Routing) 방식이다. 라우터 LLM은 현재 상태 $C_t$와 에이전트의 결과물 $O_t$를 입력받아, 다음에 호출할 에이전트 $A_{t+1}$을 결정하는 함수 $g(C_t, O_t) \rightarrow A_{t+1}$를 수행한다. 예를 들어, 검수 에이전트의 피드백이 매우 치명적이라고 판단될 경우, 수정 에이전트로 보내지 않고 다시 설계 단계(Planner)로 되돌려 구조부터 재검토하게 하는 등 유연한 경로 변경이 가능하다. 이는 작업의 복잡도에 따라 최적의 경로를 스스로 탐색하므로 품질 향상 가능성이 높으나, 라우팅을 위한 추가적인 토큰 소비와 추론 시간이 발생한다.

\subsubsection{Hybrid 오케스트레이션}
Hybrid 구조는 전체적인 큰 흐름을 제어하는 매크로 플로우(Macro-flow)와 세부 단계 내의 흐름을 제어하는 마이크로 플로우(Micro-flow)를 결합한 계층적 제어 방식이다. 
\begin{itemize}
    \item \textbf{Macro-flow (Code):} 전체 공정을 [설계] $\rightarrow$ [초안 작성 및 보완] $\rightarrow$ [최종 마무리]의 큰 단계로 나누어 결정론적으로 제어한다.
    \item \textbf{Micro-flow (LLM):} [초안 작성 및 보완] 단계 내부에서는 라우터 LLM이 집필 에이전트와 검수 에이전트 사이의 루프를 동적으로 관리한다. 즉, 검수 에이전트가 만족할 만한 품질 수준에 도달할 때까지 LLM 기반의 동적 라우팅을 통해 반복 수정 작업을 수행하며, 목표 품질에 도달하면 다시 매크로 플로우의 다음 단계로 전이한다.
\end{itemize}
이러한 계층적 구조는 Code 방식의 효율성과 LLM 방식의 유연성을 동시에 확보하여, 전체 실행 시간의 예측 가능성을 높이면서도 세부 품질을 정밀하게 제어하는 것을 목적으로 한다.